\textbf{Phần 1. Trắc nghiệm lựa chọn}
\setcounterpageref{ex}{0}

\begin{ex} %[2D10V49]
Số tròn trăm liền trước của số $600$ là
\choice
{$590$}
{\True $500$}
{$700$}
{$599$}
\loigiai{
Số tròn trăm liền trước của 600 là 500.
\textbf{Đáp án của câu hỏi là B}
}
\end{ex}

\begin{ex} %[2D10V49]
Dãy số tròn trăm nào sau đây được viết theo thứ tự từ bé đến lớn?
\choice
{$900, 800, 700$}
{$200, 400, 300$}
{$600, 400, 200$}
{\True $300, 500, 700$}
\loigiai{
Dãy $300, 500, 700$ xếp tăng dần.
\textbf{Đáp án của câu hỏi là D}
}
\end{ex}

\begin{ex} %[2D10V49]
Số gồm $7$ trăm và $0$ chục, $0$ đơn vị viết là
\choice
{\True $700$}
{$70$}
{$707$}
{$7$}
\loigiai{
Viết là $700$.
\textbf{Đáp án của câu hỏi là A}
}
\end{ex}

\begin{ex} %[2D10V49]
Từ $100$ đến $900$ có tất cả bao nhiêu số tròn trăm?
\choice
{$8$ số}
{$10$ số}
{\True $9$ số}
{$7$ số}
\loigiai{
Các số tròn trăm: 100, 200, 300, 400, 500, 600, 700, 800, 900. Có 9 số.
\textbf{Đáp án của câu hỏi là C}
}
\end{ex}

\begin{ex} %[2D10V49]
Số tròn chục lớn nhất có hai chữ số đem cộng với $100$ thì được kết quả là
\choice
{$199$}
{\True $190$}
{$180$}
{$200$}
\loigiai{
Số tròn chục lớn nhất có hai chữ số là $90$. Ta có: $90 + 100 = 190$.
\textbf{Đáp án của câu hỏi là B}
}
\end{ex}

\begin{ex} %[2D10V49]
Trong các số sau, số nào vừa là số tròn trăm vừa có chữ số hàng trăm là số chẵn lớn hơn $5$?
\choice
{$500$}
{$700$}
{$650$}
{\True $800$}
\loigiai{
Chữ số hàng trăm là số chẵn lớn hơn 5 chỉ có 6 hoặc 8. Ở đây có 800 là số tròn trăm.
\textbf{Đáp án của câu hỏi là D}
}
\end{ex}

\begin{ex} %[2D10V49]
Một nông trại đóng gói $400$ quả trứng gà, mỗi thùng xếp đúng $100$ quả. Hỏi nông trại xếp được bao nhiêu thùng trứng?
\choice
{\True $4$ thùng}
{$40$ thùng}
{$14$ thùng}
{$400$ thùng}
\loigiai{
$400 : 100 = 4$ (thùng).
\textbf{Đáp án của câu hỏi là A}
}
\end{ex}

\begin{ex} %[2D10V49]
Tổng của số tròn trăm bé nhất có ba chữ số và số tròn trăm lớn nhất có ba chữ số là
\choice
{$900$}
{$800$}
{\True $1 000$}
{$1 100$}
\loigiai{
Số tròn trăm bé nhất có 3 chữ số là 100, lớn nhất là 900. Tổng: $100 + 900 = 1 000$.
\textbf{Đáp án của câu hỏi là C}
}
\end{ex}

\begin{ex} %[2D10V49]
Số tròn trăm nào nằm chính giữa $400$ và $600$?
\choice
{$450$}
{\True $500$}
{$550$}
{$300$}
\loigiai{
Số tròn trăm ở chính giữa 400 và 600 là 500.
\textbf{Đáp án của câu hỏi là B}
}
\end{ex}

\begin{ex} %[2D10V49]
Có bao nhiêu số tròn chục lớn hơn $120$ và bé hơn $180$?
\choice
{$4$ số}
{$6$ số}
{$7$ số}
{\True $5$ số}
\loigiai{
Các số là: 130, 140, 150, 160, 170. Có tất cả 5 số.
\textbf{Đáp án của câu hỏi là D}
}
\end{ex}

\begin{ex} %[2D10V49]
Một trang trại nuôi $300$ con gà trống và $500$ con gà mái. Trang trại đó nuôi tất cả bao nhiêu con gà?
\choice
{\True $800$ con}
{$700$ con}
{$900$ con}
{$600$ con}
\loigiai{
$300 + 500 = 800$ (con gà).
\textbf{Đáp án của câu hỏi là A}
}
\end{ex}

\begin{ex} %[2D10V49]
Một bồn chứa $900$ lít nước, người ta đã xả đi $300$ lít nước để tưới vườn. Bồn còn lại bao nhiêu lít nước?
\choice
{$500$ lít}
{$400$ lít}
{\True $600$ lít}
{$700$ lít}
\loigiai{
$900 - 300 = 600$ (lít).
\textbf{Đáp án của câu hỏi là C}
}
\end{ex}

\textbf{Phần 2. Trắc nghiệm đúng sai}
\setcounterpageref{ex}{0}

\begin{ex} %[2D10V49]
Xét tính đúng sai của các khẳng định về số tròn trăm, tròn chục:
\choiceTF
{\True Số 1 000 là một số tròn trăm.}
{Số 350 là một số tròn trăm.}
{\True Giữa hai số tròn trăm 200 và 400 chỉ có đúng một số tròn trăm là 300.}
{\True Tổng của 200 và 800 bằng 1 000.}
\loigiai{
\begin{itemize}
    \item Ý a) Đúng: 1 000 gồm 10 trăm, tận cùng là hai chữ số 0 nên là số tròn trăm.
    \item Ý b) Sai: 350 là số tròn chục nhưng không phải số tròn trăm.
    \item Ý c) Đúng: Chính xác, chỉ có 300.
    \item Ý d) Đúng: 200 + 800 = 1 000.
\end{itemize}
\begin{itemize}
    \item a) \textbf{Đúng}.
    \item b) \textbf{Sai}.
    \item c) \textbf{Đúng}.
    \item d) \textbf{Đúng}.
\end{itemize}
}
\end{ex}

\begin{ex} %[2D10V49]
Cho ba số tròn trăm A = 200, B = 400, C = 600:
\choiceTF
{\True Tổng của A và B bằng C.}
{\True C gấp 3 lần A.}
{Hiệu giữa C và B bé hơn A.}
{\True Tổng của cả ba số là 1 200.}
\loigiai{
\begin{itemize}
    \item Ý a) Đúng: 200 + 400 = 600.
    \item Ý b) Đúng: 600 : 200 = 3.
    \item Ý c) Sai: C - B = 600 - 400 = 200 = A, không bé hơn.
    \item Ý d) Đúng: 200 + 400 + 600 = 1 200.
\end{itemize}
\begin{itemize}
    \item a) \textbf{Đúng}.
    \item b) \textbf{Đúng}.
    \item c) \textbf{Sai}.
    \item d) \textbf{Đúng}.
\end{itemize}
}
\end{ex}

\textbf{Phần 3. Trả lời ngắn}
\setcounterpageref{ex}{0}

\begin{ex} %[2D10V49]
Tính tổng của số tròn trăm lớn nhất có ba chữ số và số tròn chục bé nhất có hai chữ số.
\shortans{910}
\loigiai{
\textbf{Đáp số:} $910$\\
Số tròn trăm lớn nhất có 3 chữ số là 900. Số tròn chục bé nhất có 2 chữ số là 10. Tổng: $900 + 10 = 910$.
}
\end{ex}

\begin{ex} %[2D10V49]
Kho A có 400 bao gạo, kho B có ít hơn kho A 100 bao gạo. Cả hai kho có tất cả bao nhiêu bao gạo?
\shortans{700}
\loigiai{
\textbf{Đáp số:} $700$\\
Kho B có: $400 - 100 = 300$ bao. Cả hai kho: $400 + 300 = 700$ bao.
}
\end{ex}

\begin{ex} %[2D10V49]
Tìm một số tròn trăm, biết rằng nếu lấy 1 000 trừ đi số đó thì được hiệu là 600.
\shortans{400}
\loigiai{
\textbf{Đáp số:} $400$\\
Số cần tìm: $1 000 - 600 = 400$.
}
\end{ex}

\begin{ex} %[2D10V49]
Có bao nhiêu số tròn chục có hai chữ số?
\shortans{9}
\loigiai{
\textbf{Đáp số:} $9$\\
Đó là 10, 20, ..., 90. Có 9 số.
}
\end{ex}

\begin{ex} %[2D10V49]
Một cửa hàng ngày thứ nhất bán 200 m vải, ngày thứ hai bán 300 m vải. Cả hai ngày cửa hàng bán được bao nhiêu mét vải?
\shortans{500}
\loigiai{
\textbf{Đáp số:} $500$\\
200 + 300 = 500 m.
}
\end{ex}

\begin{ex} %[2D10V49]
Tìm số tròn trăm $x$ thỏa mãn: $x + 300 = 700$.
\shortans{400}
\loigiai{
\textbf{Đáp số:} $400$\\
$x = 700 - 300 = 400$.
}
\end{ex}
